\documentclass[12pt,a4paper]{article}
\usepackage[top=1.8cm,bottom=1.8cm,left=1.8cm,right=1.8cm]{geometry}
\usepackage{fontspec}
\setmainfont{TeX Gyre Termes}
\usepackage{amsmath,amssymb}
\usepackage{fancyhdr}
\usepackage{tasks}
\usepackage{tikz}
\settasks{label=\textbf{\Alph*.},label-width=1.5em,item-indent=2em}
\setlength{\parindent}{0pt}
\setlength{\parskip}{3pt}
\pagestyle{fancy}\fancyhf{}\renewcommand{\headrulewidth}{0pt}
\fancyhead[L]{\small Ngân hàng bài tập Toán 7 -- 7H2-C7-B25-D03}
\cfoot{\small Trang \thepage}
\newif\ifloigiai
\loigiaitrue
\newenvironment{cauhoi}[2]{%
  \par\medskip\noindent\rule{\textwidth}{0.4pt}\par\nopagebreak
  \noindent\textbf{[#1]}\hfill{\scriptsize\ttfamily #2}\par\nopagebreak\smallskip}{\par}
\newcommand{\loigiai}[1]{\ifloigiai\par\smallskip\noindent\textit{Lời giải.} #1\par\fi}
\newcommand{\dapan}[1]{\ifloigiai\par\smallskip\noindent\textit{Đáp án:} \textbf{#1}\par\fi}
\begin{document}
\begin{center}
{\Large\textbf{NGÂN HÀNG BÀI TẬP TOÁN 7}}\\[3pt]
\textbf{Chương 7. Biểu thức đại số và đa thức một biến}\\[2pt]
\textbf{Bài 25. Đa thức một biến}\\[2pt]
\textbf{Dạng D03: Tìm hệ số để có nghiệm cho trước}\\[2pt]
\textit{Mã dạng: 7H2-C7-B25-D03 \quad 50 bài tự luận + 50 câu trắc nghiệm \quad (VD: 50)}
\end{center}
\textbf{Quy ước mã câu:} \textit{khối, học kì, chương, bài, dạng, loại (TL/TN), mức độ, số thứ tự}.\par Ví dụ: \texttt{7H2-C7-B25-D03-TL-NB-001}.\par Để in bản cho học sinh, đặt \verb|\loigiaifalse| ngay sau \verb|\loigiaitrue|.
\section*{Phần I. Tự luận}
\begin{cauhoi}{Bài 1}{7H2-C7-B25-D03-TL-VD-001}
Một khối hình ghép gồm hai hình hộp chữ nhật xếp chồng lên nhau, cùng đáy $7$ cm $\times$ $6$ cm, chiều cao lần lượt $3$ cm và $10$ cm. Tính thể tích khối ghép.
\loigiai{Thể tích hai khối là $7\cdot 6\cdot 3=126$ và $7\cdot 6\cdot 10=420$ ($\text{cm}^3$). Thể tích khối ghép là $126+420=546$ ($\text{cm}^3$).}
\end{cauhoi}
\begin{cauhoi}{Bài 2}{7H2-C7-B25-D03-TL-VD-002}
Một lăng trụ đứng có đáy là tam giác vuông với hai cạnh góc vuông $8$ cm và $2$ cm, chiều cao lăng trụ $9$ cm. Tính thể tích.
\loigiai{Diện tích đáy là $\dfrac{8\cdot 2}{2}=8$ ($\text{cm}^2$). Thể tích là $8\cdot 9=72$ ($\text{cm}^3$).}
\end{cauhoi}
\begin{cauhoi}{Bài 3}{7H2-C7-B25-D03-TL-VD-003}
Tìm hệ số $a$ để $x=-6$ là nghiệm của đa thức $P(x)=x^{2}+ax$.
\loigiai{$x=-6$ là nghiệm nên $P(-6)=0$: $36-6\cdot a=0$ nên $-6a=-36$, suy ra $a=6$.}
\end{cauhoi}
\begin{cauhoi}{Bài 4}{7H2-C7-B25-D03-TL-VD-004}
Tìm hệ số $a$ để $x=-1$ là nghiệm của đa thức $P(x)=x^{2}+ax+2$.
\loigiai{$x=-1$ là nghiệm nên $P(-1)=0$: $1-1\cdot a+2=0$ nên $-a=-3$, suy ra $a=3$.}
\end{cauhoi}
\begin{cauhoi}{Bài 5}{7H2-C7-B25-D03-TL-VD-005}
Tìm hệ số $a$ để $x=1$ là nghiệm của đa thức $P(x)=x^{2}+ax-8$.
\loigiai{$x=1$ là nghiệm nên $P(1)=0$: $1+1\cdot a-8=0$ nên $a=7$, suy ra $a=7$.}
\end{cauhoi}
\begin{cauhoi}{Bài 6}{7H2-C7-B25-D03-TL-VD-006}
Tìm hệ số $a$ để $x=1$ là nghiệm của đa thức $P(x)=x^{2}+ax-6$.
\loigiai{$x=1$ là nghiệm nên $P(1)=0$: $1+1\cdot a-6=0$ nên $a=5$, suy ra $a=5$.}
\end{cauhoi}
\begin{cauhoi}{Bài 7}{7H2-C7-B25-D03-TL-VD-007}
Một lăng trụ đứng có đáy là tam giác vuông với hai cạnh góc vuông $8$ cm và $10$ cm, chiều cao lăng trụ $4$ cm. Tính thể tích.
\loigiai{Diện tích đáy là $\dfrac{8\cdot 10}{2}=40$ ($\text{cm}^2$). Thể tích là $40\cdot 4=160$ ($\text{cm}^3$).}
\end{cauhoi}
\begin{cauhoi}{Bài 8}{7H2-C7-B25-D03-TL-VD-008}
Một lăng trụ đứng có đáy là tam giác vuông với hai cạnh góc vuông $9$ cm và $8$ cm, chiều cao lăng trụ $11$ cm. Tính thể tích.
\loigiai{Diện tích đáy là $\dfrac{9\cdot 8}{2}=36$ ($\text{cm}^2$). Thể tích là $36\cdot 11=396$ ($\text{cm}^3$).}
\end{cauhoi}
\begin{cauhoi}{Bài 9}{7H2-C7-B25-D03-TL-VD-009}
Một khối hình ghép gồm hai hình hộp chữ nhật xếp chồng lên nhau, cùng đáy $10$ cm $\times$ $2$ cm, chiều cao lần lượt $5$ cm và $4$ cm. Tính thể tích khối ghép.
\loigiai{Thể tích hai khối là $10\cdot 2\cdot 5=100$ và $10\cdot 2\cdot 4=80$ ($\text{cm}^3$). Thể tích khối ghép là $100+80=180$ ($\text{cm}^3$).}
\end{cauhoi}
\begin{cauhoi}{Bài 10}{7H2-C7-B25-D03-TL-VD-010}
Tìm hệ số $a$ để $x=-1$ là nghiệm của đa thức $P(x)=x^{2}+ax+8$.
\loigiai{$x=-1$ là nghiệm nên $P(-1)=0$: $1-1\cdot a+8=0$ nên $-a=-9$, suy ra $a=9$.}
\end{cauhoi}
\begin{cauhoi}{Bài 11}{7H2-C7-B25-D03-TL-VD-011}
Tìm hệ số $a$ để $x=-2$ là nghiệm của đa thức $P(x)=x^{2}+ax+8$.
\loigiai{$x=-2$ là nghiệm nên $P(-2)=0$: $4-2\cdot a+8=0$ nên $-2a=-12$, suy ra $a=6$.}
\end{cauhoi}
\begin{cauhoi}{Bài 12}{7H2-C7-B25-D03-TL-VD-012}
Một lăng trụ đứng có đáy là tam giác vuông với hai cạnh góc vuông $2$ cm và $2$ cm, chiều cao lăng trụ $7$ cm. Tính thể tích.
\loigiai{Diện tích đáy là $\dfrac{2\cdot 2}{2}=2$ ($\text{cm}^2$). Thể tích là $2\cdot 7=14$ ($\text{cm}^3$).}
\end{cauhoi}
\begin{cauhoi}{Bài 13}{7H2-C7-B25-D03-TL-VD-013}
Một khối hình ghép gồm hai hình hộp chữ nhật xếp chồng lên nhau, cùng đáy $7$ cm $\times$ $10$ cm, chiều cao lần lượt $6$ cm và $2$ cm. Tính thể tích khối ghép.
\loigiai{Thể tích hai khối là $7\cdot 10\cdot 6=420$ và $7\cdot 10\cdot 2=140$ ($\text{cm}^3$). Thể tích khối ghép là $420+140=560$ ($\text{cm}^3$).}
\end{cauhoi}
\begin{cauhoi}{Bài 14}{7H2-C7-B25-D03-TL-VD-014}
Một lăng trụ đứng có đáy là tam giác vuông với hai cạnh góc vuông $2$ cm và $9$ cm, chiều cao lăng trụ $8$ cm. Tính thể tích.
\loigiai{Diện tích đáy là $\dfrac{2\cdot 9}{2}=9$ ($\text{cm}^2$). Thể tích là $9\cdot 8=72$ ($\text{cm}^3$).}
\end{cauhoi}
\begin{cauhoi}{Bài 15}{7H2-C7-B25-D03-TL-VD-015}
Một lăng trụ đứng có đáy là tam giác vuông với hai cạnh góc vuông $8$ cm và $2$ cm, chiều cao lăng trụ $12$ cm. Tính thể tích.
\loigiai{Diện tích đáy là $\dfrac{8\cdot 2}{2}=8$ ($\text{cm}^2$). Thể tích là $8\cdot 12=96$ ($\text{cm}^3$).}
\end{cauhoi}
\begin{cauhoi}{Bài 16}{7H2-C7-B25-D03-TL-VD-016}
Một lăng trụ đứng có đáy là tam giác vuông với hai cạnh góc vuông $6$ cm và $3$ cm, chiều cao lăng trụ $5$ cm. Tính thể tích.
\loigiai{Diện tích đáy là $\dfrac{6\cdot 3}{2}=9$ ($\text{cm}^2$). Thể tích là $9\cdot 5=45$ ($\text{cm}^3$).}
\end{cauhoi}
\begin{cauhoi}{Bài 17}{7H2-C7-B25-D03-TL-VD-017}
Một lăng trụ đứng có đáy là tam giác vuông với hai cạnh góc vuông $4$ cm và $7$ cm, chiều cao lăng trụ $6$ cm. Tính thể tích.
\loigiai{Diện tích đáy là $\dfrac{4\cdot 7}{2}=14$ ($\text{cm}^2$). Thể tích là $14\cdot 6=84$ ($\text{cm}^3$).}
\end{cauhoi}
\begin{cauhoi}{Bài 18}{7H2-C7-B25-D03-TL-VD-018}
Tìm hệ số $a$ để $x=-3$ là nghiệm của đa thức $P(x)=x^{2}+ax-6$.
\loigiai{$x=-3$ là nghiệm nên $P(-3)=0$: $9-3\cdot a-6=0$ nên $-3a=-3$, suy ra $a=1$.}
\end{cauhoi}
\begin{cauhoi}{Bài 19}{7H2-C7-B25-D03-TL-VD-019}
Một khối hình ghép gồm hai hình hộp chữ nhật xếp chồng lên nhau, cùng đáy $9$ cm $\times$ $10$ cm, chiều cao lần lượt $5$ cm và $5$ cm. Tính thể tích khối ghép.
\loigiai{Thể tích hai khối là $9\cdot 10\cdot 5=450$ và $9\cdot 10\cdot 5=450$ ($\text{cm}^3$). Thể tích khối ghép là $450+450=900$ ($\text{cm}^3$).}
\end{cauhoi}
\begin{cauhoi}{Bài 20}{7H2-C7-B25-D03-TL-VD-020}
Một lăng trụ đứng có đáy là tam giác vuông với hai cạnh góc vuông $7$ cm và $7$ cm, chiều cao lăng trụ $2$ cm. Tính thể tích.
\loigiai{Diện tích đáy là $\dfrac{7\cdot 7}{2}=24,5$ ($\text{cm}^2$). Thể tích là $24,5\cdot 2=49$ ($\text{cm}^3$).}
\end{cauhoi}
\begin{cauhoi}{Bài 21}{7H2-C7-B25-D03-TL-VD-021}
Tìm hệ số $a$ để $x=5$ là nghiệm của đa thức $P(x)=x^{2}+ax$.
\loigiai{$x=5$ là nghiệm nên $P(5)=0$: $25+5\cdot a=0$ nên $5a=-25$, suy ra $a=-5$.}
\end{cauhoi}
\begin{cauhoi}{Bài 22}{7H2-C7-B25-D03-TL-VD-022}
Tìm hệ số $a$ để $x=1$ là nghiệm của đa thức $P(x)=x^{2}+ax+9$.
\loigiai{$x=1$ là nghiệm nên $P(1)=0$: $1+1\cdot a+9=0$ nên $a=-10$, suy ra $a=-10$.}
\end{cauhoi}
\begin{cauhoi}{Bài 23}{7H2-C7-B25-D03-TL-VD-023}
Một lăng trụ đứng có đáy là tam giác vuông với hai cạnh góc vuông $8$ cm và $3$ cm, chiều cao lăng trụ $12$ cm. Tính thể tích.
\loigiai{Diện tích đáy là $\dfrac{8\cdot 3}{2}=12$ ($\text{cm}^2$). Thể tích là $12\cdot 12=144$ ($\text{cm}^3$).}
\end{cauhoi}
\begin{cauhoi}{Bài 24}{7H2-C7-B25-D03-TL-VD-024}
Tìm hệ số $a$ để $x=3$ là nghiệm của đa thức $P(x)=x^{2}+ax-9$.
\loigiai{$x=3$ là nghiệm nên $P(3)=0$: $9+3\cdot a-9=0$ nên $3a=0$, suy ra $a=0$.}
\end{cauhoi}
\begin{cauhoi}{Bài 25}{7H2-C7-B25-D03-TL-VD-025}
Tìm hệ số $a$ để $x=2$ là nghiệm của đa thức $P(x)=x^{2}+ax-8$.
\loigiai{$x=2$ là nghiệm nên $P(2)=0$: $4+2\cdot a-8=0$ nên $2a=4$, suy ra $a=2$.}
\end{cauhoi}
\begin{cauhoi}{Bài 26}{7H2-C7-B25-D03-TL-VD-026}
Tìm hệ số $a$ để $x=4$ là nghiệm của đa thức $P(x)=x^{2}+ax-4$.
\loigiai{$x=4$ là nghiệm nên $P(4)=0$: $16+4\cdot a-4=0$ nên $4a=-12$, suy ra $a=-3$.}
\end{cauhoi}
\begin{cauhoi}{Bài 27}{7H2-C7-B25-D03-TL-VD-027}
Tìm hệ số $a$ để $x=4$ là nghiệm của đa thức $P(x)=x^{2}+ax$.
\loigiai{$x=4$ là nghiệm nên $P(4)=0$: $16+4\cdot a=0$ nên $4a=-16$, suy ra $a=-4$.}
\end{cauhoi}
\begin{cauhoi}{Bài 28}{7H2-C7-B25-D03-TL-VD-028}
Tìm hệ số $a$ để $x=-3$ là nghiệm của đa thức $P(x)=x^{2}+ax+9$.
\loigiai{$x=-3$ là nghiệm nên $P(-3)=0$: $9-3\cdot a+9=0$ nên $-3a=-18$, suy ra $a=6$.}
\end{cauhoi}
\begin{cauhoi}{Bài 29}{7H2-C7-B25-D03-TL-VD-029}
Tìm hệ số $a$ để $x=-2$ là nghiệm của đa thức $P(x)=x^{2}+ax-8$.
\loigiai{$x=-2$ là nghiệm nên $P(-2)=0$: $4-2\cdot a-8=0$ nên $-2a=4$, suy ra $a=-2$.}
\end{cauhoi}
\begin{cauhoi}{Bài 30}{7H2-C7-B25-D03-TL-VD-030}
Một khối hình ghép gồm hai hình hộp chữ nhật xếp chồng lên nhau, cùng đáy $6$ cm $\times$ $2$ cm, chiều cao lần lượt $7$ cm và $2$ cm. Tính thể tích khối ghép.
\loigiai{Thể tích hai khối là $6\cdot 2\cdot 7=84$ và $6\cdot 2\cdot 2=24$ ($\text{cm}^3$). Thể tích khối ghép là $84+24=108$ ($\text{cm}^3$).}
\end{cauhoi}
\begin{cauhoi}{Bài 31}{7H2-C7-B25-D03-TL-VD-031}
Một khối hình ghép gồm hai hình hộp chữ nhật xếp chồng lên nhau, cùng đáy $4$ cm $\times$ $8$ cm, chiều cao lần lượt $8$ cm và $6$ cm. Tính thể tích khối ghép.
\loigiai{Thể tích hai khối là $4\cdot 8\cdot 8=256$ và $4\cdot 8\cdot 6=192$ ($\text{cm}^3$). Thể tích khối ghép là $256+192=448$ ($\text{cm}^3$).}
\end{cauhoi}
\begin{cauhoi}{Bài 32}{7H2-C7-B25-D03-TL-VD-032}
Tìm hệ số $a$ để $x=1$ là nghiệm của đa thức $P(x)=x^{2}+ax$.
\loigiai{$x=1$ là nghiệm nên $P(1)=0$: $1+1\cdot a=0$ nên $a=-1$, suy ra $a=-1$.}
\end{cauhoi}
\begin{cauhoi}{Bài 33}{7H2-C7-B25-D03-TL-VD-033}
Tìm hệ số $a$ để $x=1$ là nghiệm của đa thức $P(x)=x^{2}+ax-3$.
\loigiai{$x=1$ là nghiệm nên $P(1)=0$: $1+1\cdot a-3=0$ nên $a=2$, suy ra $a=2$.}
\end{cauhoi}
\begin{cauhoi}{Bài 34}{7H2-C7-B25-D03-TL-VD-034}
Một lăng trụ đứng có đáy là tam giác vuông với hai cạnh góc vuông $4$ cm và $5$ cm, chiều cao lăng trụ $10$ cm. Tính thể tích.
\loigiai{Diện tích đáy là $\dfrac{4\cdot 5}{2}=10$ ($\text{cm}^2$). Thể tích là $10\cdot 10=100$ ($\text{cm}^3$).}
\end{cauhoi}
\begin{cauhoi}{Bài 35}{7H2-C7-B25-D03-TL-VD-035}
Một khối hình ghép gồm hai hình hộp chữ nhật xếp chồng lên nhau, cùng đáy $5$ cm $\times$ $2$ cm, chiều cao lần lượt $7$ cm và $6$ cm. Tính thể tích khối ghép.
\loigiai{Thể tích hai khối là $5\cdot 2\cdot 7=70$ và $5\cdot 2\cdot 6=60$ ($\text{cm}^3$). Thể tích khối ghép là $70+60=130$ ($\text{cm}^3$).}
\end{cauhoi}
\begin{cauhoi}{Bài 36}{7H2-C7-B25-D03-TL-VD-036}
Tìm hệ số $a$ để $x=4$ là nghiệm của đa thức $P(x)=x^{2}+ax-8$.
\loigiai{$x=4$ là nghiệm nên $P(4)=0$: $16+4\cdot a-8=0$ nên $4a=-8$, suy ra $a=-2$.}
\end{cauhoi}
\begin{cauhoi}{Bài 37}{7H2-C7-B25-D03-TL-VD-037}
Một lăng trụ đứng có đáy là tam giác vuông với hai cạnh góc vuông $5$ cm và $5$ cm, chiều cao lăng trụ $11$ cm. Tính thể tích.
\loigiai{Diện tích đáy là $\dfrac{5\cdot 5}{2}=12,5$ ($\text{cm}^2$). Thể tích là $12,5\cdot 11=137,5$ ($\text{cm}^3$).}
\end{cauhoi}
\begin{cauhoi}{Bài 38}{7H2-C7-B25-D03-TL-VD-038}
Một khối hình ghép gồm hai hình hộp chữ nhật xếp chồng lên nhau, cùng đáy $3$ cm $\times$ $3$ cm, chiều cao lần lượt $7$ cm và $4$ cm. Tính thể tích khối ghép.
\loigiai{Thể tích hai khối là $3\cdot 3\cdot 7=63$ và $3\cdot 3\cdot 4=36$ ($\text{cm}^3$). Thể tích khối ghép là $63+36=99$ ($\text{cm}^3$).}
\end{cauhoi}
\begin{cauhoi}{Bài 39}{7H2-C7-B25-D03-TL-VD-039}
Một khối hình ghép gồm hai hình hộp chữ nhật xếp chồng lên nhau, cùng đáy $7$ cm $\times$ $5$ cm, chiều cao lần lượt $10$ cm và $8$ cm. Tính thể tích khối ghép.
\loigiai{Thể tích hai khối là $7\cdot 5\cdot 10=350$ và $7\cdot 5\cdot 8=280$ ($\text{cm}^3$). Thể tích khối ghép là $350+280=630$ ($\text{cm}^3$).}
\end{cauhoi}
\begin{cauhoi}{Bài 40}{7H2-C7-B25-D03-TL-VD-040}
Một khối hình ghép gồm hai hình hộp chữ nhật xếp chồng lên nhau, cùng đáy $3$ cm $\times$ $7$ cm, chiều cao lần lượt $2$ cm và $8$ cm. Tính thể tích khối ghép.
\loigiai{Thể tích hai khối là $3\cdot 7\cdot 2=42$ và $3\cdot 7\cdot 8=168$ ($\text{cm}^3$). Thể tích khối ghép là $42+168=210$ ($\text{cm}^3$).}
\end{cauhoi}
\begin{cauhoi}{Bài 41}{7H2-C7-B25-D03-TL-VD-041}
Tìm hệ số $a$ để $x=3$ là nghiệm của đa thức $P(x)=x^{2}+ax+9$.
\loigiai{$x=3$ là nghiệm nên $P(3)=0$: $9+3\cdot a+9=0$ nên $3a=-18$, suy ra $a=-6$.}
\end{cauhoi}
\begin{cauhoi}{Bài 42}{7H2-C7-B25-D03-TL-VD-042}
Một lăng trụ đứng có đáy là tam giác vuông với hai cạnh góc vuông $7$ cm và $3$ cm, chiều cao lăng trụ $11$ cm. Tính thể tích.
\loigiai{Diện tích đáy là $\dfrac{7\cdot 3}{2}=10,5$ ($\text{cm}^2$). Thể tích là $10,5\cdot 11=115,5$ ($\text{cm}^3$).}
\end{cauhoi}
\begin{cauhoi}{Bài 43}{7H2-C7-B25-D03-TL-VD-043}
Tìm hệ số $a$ để $x=-1$ là nghiệm của đa thức $P(x)=x^{2}+ax-4$.
\loigiai{$x=-1$ là nghiệm nên $P(-1)=0$: $1-1\cdot a-4=0$ nên $-a=3$, suy ra $a=-3$.}
\end{cauhoi}
\begin{cauhoi}{Bài 44}{7H2-C7-B25-D03-TL-VD-044}
Một khối hình ghép gồm hai hình hộp chữ nhật xếp chồng lên nhau, cùng đáy $3$ cm $\times$ $3$ cm, chiều cao lần lượt $7$ cm và $8$ cm. Tính thể tích khối ghép.
\loigiai{Thể tích hai khối là $3\cdot 3\cdot 7=63$ và $3\cdot 3\cdot 8=72$ ($\text{cm}^3$). Thể tích khối ghép là $63+72=135$ ($\text{cm}^3$).}
\end{cauhoi}
\begin{cauhoi}{Bài 45}{7H2-C7-B25-D03-TL-VD-045}
Một khối hình ghép gồm hai hình hộp chữ nhật xếp chồng lên nhau, cùng đáy $7$ cm $\times$ $4$ cm, chiều cao lần lượt $9$ cm và $10$ cm. Tính thể tích khối ghép.
\loigiai{Thể tích hai khối là $7\cdot 4\cdot 9=252$ và $7\cdot 4\cdot 10=280$ ($\text{cm}^3$). Thể tích khối ghép là $252+280=532$ ($\text{cm}^3$).}
\end{cauhoi}
\begin{cauhoi}{Bài 46}{7H2-C7-B25-D03-TL-VD-046}
Một khối hình ghép gồm hai hình hộp chữ nhật xếp chồng lên nhau, cùng đáy $7$ cm $\times$ $3$ cm, chiều cao lần lượt $4$ cm và $6$ cm. Tính thể tích khối ghép.
\loigiai{Thể tích hai khối là $7\cdot 3\cdot 4=84$ và $7\cdot 3\cdot 6=126$ ($\text{cm}^3$). Thể tích khối ghép là $84+126=210$ ($\text{cm}^3$).}
\end{cauhoi}
\begin{cauhoi}{Bài 47}{7H2-C7-B25-D03-TL-VD-047}
Tìm hệ số $a$ để $x=-3$ là nghiệm của đa thức $P(x)=x^{2}+ax+6$.
\loigiai{$x=-3$ là nghiệm nên $P(-3)=0$: $9-3\cdot a+6=0$ nên $-3a=-15$, suy ra $a=5$.}
\end{cauhoi}
\begin{cauhoi}{Bài 48}{7H2-C7-B25-D03-TL-VD-048}
Tìm hệ số $a$ để $x=1$ là nghiệm của đa thức $P(x)=x^{2}+ax-5$.
\loigiai{$x=1$ là nghiệm nên $P(1)=0$: $1+1\cdot a-5=0$ nên $a=4$, suy ra $a=4$.}
\end{cauhoi}
\begin{cauhoi}{Bài 49}{7H2-C7-B25-D03-TL-VD-049}
Một lăng trụ đứng có đáy là tam giác vuông với hai cạnh góc vuông $9$ cm và $10$ cm, chiều cao lăng trụ $7$ cm. Tính thể tích.
\loigiai{Diện tích đáy là $\dfrac{9\cdot 10}{2}=45$ ($\text{cm}^2$). Thể tích là $45\cdot 7=315$ ($\text{cm}^3$).}
\end{cauhoi}
\begin{cauhoi}{Bài 50}{7H2-C7-B25-D03-TL-VD-050}
Một lăng trụ đứng có đáy là tam giác vuông với hai cạnh góc vuông $8$ cm và $6$ cm, chiều cao lăng trụ $4$ cm. Tính thể tích.
\loigiai{Diện tích đáy là $\dfrac{8\cdot 6}{2}=24$ ($\text{cm}^2$). Thể tích là $24\cdot 4=96$ ($\text{cm}^3$).}
\end{cauhoi}
\clearpage\section*{Phần II. Trắc nghiệm}
\begin{cauhoi}{Câu 1}{7H2-C7-B25-D03-TN-VD-001}
Một khối hình ghép gồm hai hình hộp chữ nhật xếp chồng lên nhau, cùng đáy $7$ cm $\times$ $6$ cm, chiều cao lần lượt $3$ cm và $10$ cm. Tính thể tích khối ghép.
\begin{tasks}(2)
\task $420\text{ cm}^3$
\task $126\text{ cm}^3$
\task $547\text{ cm}^3$
\task $546\text{ cm}^3$
\end{tasks}
\dapan{D}
\end{cauhoi}
\begin{cauhoi}{Câu 2}{7H2-C7-B25-D03-TN-VD-002}
Một lăng trụ đứng có đáy là tam giác vuông với hai cạnh góc vuông $8$ cm và $2$ cm, chiều cao lăng trụ $9$ cm. Tính thể tích.
\begin{tasks}(2)
\task $72\text{ cm}^3$
\task $90\text{ cm}^3$
\task $144\text{ cm}^3$
\task $82\text{ cm}^3$
\end{tasks}
\dapan{A}
\end{cauhoi}
\begin{cauhoi}{Câu 3}{7H2-C7-B25-D03-TN-VD-003}
Tìm hệ số $a$ để $x=-6$ là nghiệm của đa thức $P(x)=x^{2}+ax$.
\begin{tasks}(2)
\task $-6$
\task $7$
\task $5$
\task $6$
\end{tasks}
\dapan{D}
\end{cauhoi}
\begin{cauhoi}{Câu 4}{7H2-C7-B25-D03-TN-VD-004}
Tìm hệ số $a$ để $x=-1$ là nghiệm của đa thức $P(x)=x^{2}+ax+2$.
\begin{tasks}(2)
\task $-3$
\task $3$
\task $2$
\task $4$
\end{tasks}
\dapan{B}
\end{cauhoi}
\begin{cauhoi}{Câu 5}{7H2-C7-B25-D03-TN-VD-005}
Tìm hệ số $a$ để $x=1$ là nghiệm của đa thức $P(x)=x^{2}+ax-8$.
\begin{tasks}(2)
\task $-7$
\task $8$
\task $7$
\task $6$
\end{tasks}
\dapan{C}
\end{cauhoi}
\begin{cauhoi}{Câu 6}{7H2-C7-B25-D03-TN-VD-006}
Tìm hệ số $a$ để $x=1$ là nghiệm của đa thức $P(x)=x^{2}+ax-6$.
\begin{tasks}(2)
\task $5$
\task $4$
\task $6$
\task $-5$
\end{tasks}
\dapan{A}
\end{cauhoi}
\begin{cauhoi}{Câu 7}{7H2-C7-B25-D03-TN-VD-007}
Một lăng trụ đứng có đáy là tam giác vuông với hai cạnh góc vuông $8$ cm và $10$ cm, chiều cao lăng trụ $4$ cm. Tính thể tích.
\begin{tasks}(2)
\task $320\text{ cm}^3$
\task $160\text{ cm}^3$
\task $72\text{ cm}^3$
\task $170\text{ cm}^3$
\end{tasks}
\dapan{B}
\end{cauhoi}
\begin{cauhoi}{Câu 8}{7H2-C7-B25-D03-TN-VD-008}
Một lăng trụ đứng có đáy là tam giác vuông với hai cạnh góc vuông $9$ cm và $8$ cm, chiều cao lăng trụ $11$ cm. Tính thể tích.
\begin{tasks}(2)
\task $406\text{ cm}^3$
\task $792\text{ cm}^3$
\task $396\text{ cm}^3$
\task $187\text{ cm}^3$
\end{tasks}
\dapan{C}
\end{cauhoi}
\begin{cauhoi}{Câu 9}{7H2-C7-B25-D03-TN-VD-009}
Một khối hình ghép gồm hai hình hộp chữ nhật xếp chồng lên nhau, cùng đáy $10$ cm $\times$ $2$ cm, chiều cao lần lượt $5$ cm và $4$ cm. Tính thể tích khối ghép.
\begin{tasks}(2)
\task $80\text{ cm}^3$
\task $180\text{ cm}^3$
\task $100\text{ cm}^3$
\task $181\text{ cm}^3$
\end{tasks}
\dapan{B}
\end{cauhoi}
\begin{cauhoi}{Câu 10}{7H2-C7-B25-D03-TN-VD-010}
Tìm hệ số $a$ để $x=-1$ là nghiệm của đa thức $P(x)=x^{2}+ax+8$.
\begin{tasks}(2)
\task $8$
\task $10$
\task $-9$
\task $9$
\end{tasks}
\dapan{D}
\end{cauhoi}
\begin{cauhoi}{Câu 11}{7H2-C7-B25-D03-TN-VD-011}
Tìm hệ số $a$ để $x=-2$ là nghiệm của đa thức $P(x)=x^{2}+ax+8$.
\begin{tasks}(2)
\task $7$
\task $5$
\task $-6$
\task $6$
\end{tasks}
\dapan{D}
\end{cauhoi}
\begin{cauhoi}{Câu 12}{7H2-C7-B25-D03-TN-VD-012}
Một lăng trụ đứng có đáy là tam giác vuông với hai cạnh góc vuông $2$ cm và $2$ cm, chiều cao lăng trụ $7$ cm. Tính thể tích.
\begin{tasks}(2)
\task $24\text{ cm}^3$
\task $140\text{ cm}^3$
\task $28\text{ cm}^3$
\task $14\text{ cm}^3$
\end{tasks}
\dapan{D}
\end{cauhoi}
\begin{cauhoi}{Câu 13}{7H2-C7-B25-D03-TN-VD-013}
Một khối hình ghép gồm hai hình hộp chữ nhật xếp chồng lên nhau, cùng đáy $7$ cm $\times$ $10$ cm, chiều cao lần lượt $6$ cm và $2$ cm. Tính thể tích khối ghép.
\begin{tasks}(2)
\task $420\text{ cm}^3$
\task $140\text{ cm}^3$
\task $561\text{ cm}^3$
\task $560\text{ cm}^3$
\end{tasks}
\dapan{D}
\end{cauhoi}
\begin{cauhoi}{Câu 14}{7H2-C7-B25-D03-TN-VD-014}
Một lăng trụ đứng có đáy là tam giác vuông với hai cạnh góc vuông $2$ cm và $9$ cm, chiều cao lăng trụ $8$ cm. Tính thể tích.
\begin{tasks}(2)
\task $144\text{ cm}^3$
\task $72\text{ cm}^3$
\task $82\text{ cm}^3$
\task $88\text{ cm}^3$
\end{tasks}
\dapan{B}
\end{cauhoi}
\begin{cauhoi}{Câu 15}{7H2-C7-B25-D03-TN-VD-015}
Một lăng trụ đứng có đáy là tam giác vuông với hai cạnh góc vuông $8$ cm và $2$ cm, chiều cao lăng trụ $12$ cm. Tính thể tích.
\begin{tasks}(2)
\task $96\text{ cm}^3$
\task $120\text{ cm}^3$
\task $192\text{ cm}^3$
\task $106\text{ cm}^3$
\end{tasks}
\dapan{A}
\end{cauhoi}
\begin{cauhoi}{Câu 16}{7H2-C7-B25-D03-TN-VD-016}
Một lăng trụ đứng có đáy là tam giác vuông với hai cạnh góc vuông $6$ cm và $3$ cm, chiều cao lăng trụ $5$ cm. Tính thể tích.
\begin{tasks}(2)
\task $45\text{ cm}^3$
\task $55\text{ cm}^3$
\task $450\text{ cm}^3$
\task $90\text{ cm}^3$
\end{tasks}
\dapan{A}
\end{cauhoi}
\begin{cauhoi}{Câu 17}{7H2-C7-B25-D03-TN-VD-017}
Một lăng trụ đứng có đáy là tam giác vuông với hai cạnh góc vuông $4$ cm và $7$ cm, chiều cao lăng trụ $6$ cm. Tính thể tích.
\begin{tasks}(2)
\task $94\text{ cm}^3$
\task $66\text{ cm}^3$
\task $168\text{ cm}^3$
\task $84\text{ cm}^3$
\end{tasks}
\dapan{D}
\end{cauhoi}
\begin{cauhoi}{Câu 18}{7H2-C7-B25-D03-TN-VD-018}
Tìm hệ số $a$ để $x=-3$ là nghiệm của đa thức $P(x)=x^{2}+ax-6$.
\begin{tasks}(2)
\task $0$
\task $-1$
\task $1$
\task $2$
\end{tasks}
\dapan{C}
\end{cauhoi}
\begin{cauhoi}{Câu 19}{7H2-C7-B25-D03-TN-VD-019}
Một khối hình ghép gồm hai hình hộp chữ nhật xếp chồng lên nhau, cùng đáy $9$ cm $\times$ $10$ cm, chiều cao lần lượt $5$ cm và $5$ cm. Tính thể tích khối ghép.
\begin{tasks}(2)
\task $900\text{ cm}^3$
\task $450\text{ cm}^3$
\task $901\text{ cm}^3$
\task $899\text{ cm}^3$
\end{tasks}
\dapan{A}
\end{cauhoi}
\begin{cauhoi}{Câu 20}{7H2-C7-B25-D03-TN-VD-020}
Một lăng trụ đứng có đáy là tam giác vuông với hai cạnh góc vuông $7$ cm và $7$ cm, chiều cao lăng trụ $2$ cm. Tính thể tích.
\begin{tasks}(2)
\task $59\text{ cm}^3$
\task $98\text{ cm}^3$
\task $28\text{ cm}^3$
\task $49\text{ cm}^3$
\end{tasks}
\dapan{D}
\end{cauhoi}
\begin{cauhoi}{Câu 21}{7H2-C7-B25-D03-TN-VD-021}
Tìm hệ số $a$ để $x=5$ là nghiệm của đa thức $P(x)=x^{2}+ax$.
\begin{tasks}(2)
\task $-4$
\task $5$
\task $-6$
\task $-5$
\end{tasks}
\dapan{D}
\end{cauhoi}
\begin{cauhoi}{Câu 22}{7H2-C7-B25-D03-TN-VD-022}
Tìm hệ số $a$ để $x=1$ là nghiệm của đa thức $P(x)=x^{2}+ax+9$.
\begin{tasks}(2)
\task $-9$
\task $-11$
\task $10$
\task $-10$
\end{tasks}
\dapan{D}
\end{cauhoi}
\begin{cauhoi}{Câu 23}{7H2-C7-B25-D03-TN-VD-023}
Một lăng trụ đứng có đáy là tam giác vuông với hai cạnh góc vuông $8$ cm và $3$ cm, chiều cao lăng trụ $12$ cm. Tính thể tích.
\begin{tasks}(2)
\task $132\text{ cm}^3$
\task $154\text{ cm}^3$
\task $288\text{ cm}^3$
\task $144\text{ cm}^3$
\end{tasks}
\dapan{D}
\end{cauhoi}
\begin{cauhoi}{Câu 24}{7H2-C7-B25-D03-TN-VD-024}
Tìm hệ số $a$ để $x=3$ là nghiệm của đa thức $P(x)=x^{2}+ax-9$.
\begin{tasks}(2)
\task $2$
\task $0$
\task $1$
\task $-1$
\end{tasks}
\dapan{B}
\end{cauhoi}
\begin{cauhoi}{Câu 25}{7H2-C7-B25-D03-TN-VD-025}
Tìm hệ số $a$ để $x=2$ là nghiệm của đa thức $P(x)=x^{2}+ax-8$.
\begin{tasks}(2)
\task $1$
\task $3$
\task $-2$
\task $2$
\end{tasks}
\dapan{D}
\end{cauhoi}
\begin{cauhoi}{Câu 26}{7H2-C7-B25-D03-TN-VD-026}
Tìm hệ số $a$ để $x=4$ là nghiệm của đa thức $P(x)=x^{2}+ax-4$.
\begin{tasks}(2)
\task $-2$
\task $3$
\task $-4$
\task $-3$
\end{tasks}
\dapan{D}
\end{cauhoi}
\begin{cauhoi}{Câu 27}{7H2-C7-B25-D03-TN-VD-027}
Tìm hệ số $a$ để $x=4$ là nghiệm của đa thức $P(x)=x^{2}+ax$.
\begin{tasks}(2)
\task $-5$
\task $4$
\task $-4$
\task $-3$
\end{tasks}
\dapan{C}
\end{cauhoi}
\begin{cauhoi}{Câu 28}{7H2-C7-B25-D03-TN-VD-028}
Tìm hệ số $a$ để $x=-3$ là nghiệm của đa thức $P(x)=x^{2}+ax+9$.
\begin{tasks}(2)
\task $-6$
\task $5$
\task $7$
\task $6$
\end{tasks}
\dapan{D}
\end{cauhoi}
\begin{cauhoi}{Câu 29}{7H2-C7-B25-D03-TN-VD-029}
Tìm hệ số $a$ để $x=-2$ là nghiệm của đa thức $P(x)=x^{2}+ax-8$.
\begin{tasks}(2)
\task $-1$
\task $2$
\task $-3$
\task $-2$
\end{tasks}
\dapan{D}
\end{cauhoi}
\begin{cauhoi}{Câu 30}{7H2-C7-B25-D03-TN-VD-030}
Một khối hình ghép gồm hai hình hộp chữ nhật xếp chồng lên nhau, cùng đáy $6$ cm $\times$ $2$ cm, chiều cao lần lượt $7$ cm và $2$ cm. Tính thể tích khối ghép.
\begin{tasks}(2)
\task $24\text{ cm}^3$
\task $108\text{ cm}^3$
\task $109\text{ cm}^3$
\task $84\text{ cm}^3$
\end{tasks}
\dapan{B}
\end{cauhoi}
\begin{cauhoi}{Câu 31}{7H2-C7-B25-D03-TN-VD-031}
Một khối hình ghép gồm hai hình hộp chữ nhật xếp chồng lên nhau, cùng đáy $4$ cm $\times$ $8$ cm, chiều cao lần lượt $8$ cm và $6$ cm. Tính thể tích khối ghép.
\begin{tasks}(2)
\task $256\text{ cm}^3$
\task $449\text{ cm}^3$
\task $448\text{ cm}^3$
\task $192\text{ cm}^3$
\end{tasks}
\dapan{C}
\end{cauhoi}
\begin{cauhoi}{Câu 32}{7H2-C7-B25-D03-TN-VD-032}
Tìm hệ số $a$ để $x=1$ là nghiệm của đa thức $P(x)=x^{2}+ax$.
\begin{tasks}(2)
\task $0$
\task $-2$
\task $1$
\task $-1$
\end{tasks}
\dapan{D}
\end{cauhoi}
\begin{cauhoi}{Câu 33}{7H2-C7-B25-D03-TN-VD-033}
Tìm hệ số $a$ để $x=1$ là nghiệm của đa thức $P(x)=x^{2}+ax-3$.
\begin{tasks}(2)
\task $-2$
\task $1$
\task $3$
\task $2$
\end{tasks}
\dapan{D}
\end{cauhoi}
\begin{cauhoi}{Câu 34}{7H2-C7-B25-D03-TN-VD-034}
Một lăng trụ đứng có đáy là tam giác vuông với hai cạnh góc vuông $4$ cm và $5$ cm, chiều cao lăng trụ $10$ cm. Tính thể tích.
\begin{tasks}(2)
\task $200\text{ cm}^3$
\task $100\text{ cm}^3$
\task $110\text{ cm}^3$
\task $90\text{ cm}^3$
\end{tasks}
\dapan{B}
\end{cauhoi}
\begin{cauhoi}{Câu 35}{7H2-C7-B25-D03-TN-VD-035}
Một khối hình ghép gồm hai hình hộp chữ nhật xếp chồng lên nhau, cùng đáy $5$ cm $\times$ $2$ cm, chiều cao lần lượt $7$ cm và $6$ cm. Tính thể tích khối ghép.
\begin{tasks}(2)
\task $130\text{ cm}^3$
\task $60\text{ cm}^3$
\task $70\text{ cm}^3$
\task $131\text{ cm}^3$
\end{tasks}
\dapan{A}
\end{cauhoi}
\begin{cauhoi}{Câu 36}{7H2-C7-B25-D03-TN-VD-036}
Tìm hệ số $a$ để $x=4$ là nghiệm của đa thức $P(x)=x^{2}+ax-8$.
\begin{tasks}(2)
\task $-3$
\task $2$
\task $-2$
\task $-1$
\end{tasks}
\dapan{C}
\end{cauhoi}
\begin{cauhoi}{Câu 37}{7H2-C7-B25-D03-TN-VD-037}
Một lăng trụ đứng có đáy là tam giác vuông với hai cạnh góc vuông $5$ cm và $5$ cm, chiều cao lăng trụ $11$ cm. Tính thể tích.
\begin{tasks}(2)
\task $110\text{ cm}^3$
\task $147,5\text{ cm}^3$
\task $137,5\text{ cm}^3$
\task $275\text{ cm}^3$
\end{tasks}
\dapan{C}
\end{cauhoi}
\begin{cauhoi}{Câu 38}{7H2-C7-B25-D03-TN-VD-038}
Một khối hình ghép gồm hai hình hộp chữ nhật xếp chồng lên nhau, cùng đáy $3$ cm $\times$ $3$ cm, chiều cao lần lượt $7$ cm và $4$ cm. Tính thể tích khối ghép.
\begin{tasks}(2)
\task $36\text{ cm}^3$
\task $99\text{ cm}^3$
\task $63\text{ cm}^3$
\task $100\text{ cm}^3$
\end{tasks}
\dapan{B}
\end{cauhoi}
\begin{cauhoi}{Câu 39}{7H2-C7-B25-D03-TN-VD-039}
Một khối hình ghép gồm hai hình hộp chữ nhật xếp chồng lên nhau, cùng đáy $7$ cm $\times$ $5$ cm, chiều cao lần lượt $10$ cm và $8$ cm. Tính thể tích khối ghép.
\begin{tasks}(2)
\task $631\text{ cm}^3$
\task $280\text{ cm}^3$
\task $630\text{ cm}^3$
\task $350\text{ cm}^3$
\end{tasks}
\dapan{C}
\end{cauhoi}
\begin{cauhoi}{Câu 40}{7H2-C7-B25-D03-TN-VD-040}
Một khối hình ghép gồm hai hình hộp chữ nhật xếp chồng lên nhau, cùng đáy $3$ cm $\times$ $7$ cm, chiều cao lần lượt $2$ cm và $8$ cm. Tính thể tích khối ghép.
\begin{tasks}(2)
\task $168\text{ cm}^3$
\task $211\text{ cm}^3$
\task $210\text{ cm}^3$
\task $42\text{ cm}^3$
\end{tasks}
\dapan{C}
\end{cauhoi}
\begin{cauhoi}{Câu 41}{7H2-C7-B25-D03-TN-VD-041}
Tìm hệ số $a$ để $x=3$ là nghiệm của đa thức $P(x)=x^{2}+ax+9$.
\begin{tasks}(2)
\task $-6$
\task $-7$
\task $-5$
\task $6$
\end{tasks}
\dapan{A}
\end{cauhoi}
\begin{cauhoi}{Câu 42}{7H2-C7-B25-D03-TN-VD-042}
Một lăng trụ đứng có đáy là tam giác vuông với hai cạnh góc vuông $7$ cm và $3$ cm, chiều cao lăng trụ $11$ cm. Tính thể tích.
\begin{tasks}(2)
\task $115,5\text{ cm}^3$
\task $110\text{ cm}^3$
\task $125,5\text{ cm}^3$
\task $231\text{ cm}^3$
\end{tasks}
\dapan{A}
\end{cauhoi}
\begin{cauhoi}{Câu 43}{7H2-C7-B25-D03-TN-VD-043}
Tìm hệ số $a$ để $x=-1$ là nghiệm của đa thức $P(x)=x^{2}+ax-4$.
\begin{tasks}(2)
\task $3$
\task $-2$
\task $-3$
\task $-4$
\end{tasks}
\dapan{C}
\end{cauhoi}
\begin{cauhoi}{Câu 44}{7H2-C7-B25-D03-TN-VD-044}
Một khối hình ghép gồm hai hình hộp chữ nhật xếp chồng lên nhau, cùng đáy $3$ cm $\times$ $3$ cm, chiều cao lần lượt $7$ cm và $8$ cm. Tính thể tích khối ghép.
\begin{tasks}(2)
\task $135\text{ cm}^3$
\task $72\text{ cm}^3$
\task $63\text{ cm}^3$
\task $136\text{ cm}^3$
\end{tasks}
\dapan{A}
\end{cauhoi}
\begin{cauhoi}{Câu 45}{7H2-C7-B25-D03-TN-VD-045}
Một khối hình ghép gồm hai hình hộp chữ nhật xếp chồng lên nhau, cùng đáy $7$ cm $\times$ $4$ cm, chiều cao lần lượt $9$ cm và $10$ cm. Tính thể tích khối ghép.
\begin{tasks}(2)
\task $280\text{ cm}^3$
\task $533\text{ cm}^3$
\task $252\text{ cm}^3$
\task $532\text{ cm}^3$
\end{tasks}
\dapan{D}
\end{cauhoi}
\begin{cauhoi}{Câu 46}{7H2-C7-B25-D03-TN-VD-046}
Một khối hình ghép gồm hai hình hộp chữ nhật xếp chồng lên nhau, cùng đáy $7$ cm $\times$ $3$ cm, chiều cao lần lượt $4$ cm và $6$ cm. Tính thể tích khối ghép.
\begin{tasks}(2)
\task $126\text{ cm}^3$
\task $210\text{ cm}^3$
\task $211\text{ cm}^3$
\task $84\text{ cm}^3$
\end{tasks}
\dapan{B}
\end{cauhoi}
\begin{cauhoi}{Câu 47}{7H2-C7-B25-D03-TN-VD-047}
Tìm hệ số $a$ để $x=-3$ là nghiệm của đa thức $P(x)=x^{2}+ax+6$.
\begin{tasks}(2)
\task $4$
\task $6$
\task $5$
\task $-5$
\end{tasks}
\dapan{C}
\end{cauhoi}
\begin{cauhoi}{Câu 48}{7H2-C7-B25-D03-TN-VD-048}
Tìm hệ số $a$ để $x=1$ là nghiệm của đa thức $P(x)=x^{2}+ax-5$.
\begin{tasks}(2)
\task $4$
\task $-4$
\task $3$
\task $5$
\end{tasks}
\dapan{A}
\end{cauhoi}
\begin{cauhoi}{Câu 49}{7H2-C7-B25-D03-TN-VD-049}
Một lăng trụ đứng có đáy là tam giác vuông với hai cạnh góc vuông $9$ cm và $10$ cm, chiều cao lăng trụ $7$ cm. Tính thể tích.
\begin{tasks}(2)
\task $133\text{ cm}^3$
\task $325\text{ cm}^3$
\task $315\text{ cm}^3$
\task $630\text{ cm}^3$
\end{tasks}
\dapan{C}
\end{cauhoi}
\begin{cauhoi}{Câu 50}{7H2-C7-B25-D03-TN-VD-050}
Một lăng trụ đứng có đáy là tam giác vuông với hai cạnh góc vuông $8$ cm và $6$ cm, chiều cao lăng trụ $4$ cm. Tính thể tích.
\begin{tasks}(2)
\task $106\text{ cm}^3$
\task $192\text{ cm}^3$
\task $96\text{ cm}^3$
\task $56\text{ cm}^3$
\end{tasks}
\dapan{C}
\end{cauhoi}
\clearpage\section*{Bảng đáp án trắc nghiệm}
\begin{center}\begin{tabular}{|c|c|c|c|c|c|c|c|c|c|}
\hline
1. D & 2. A & 3. D & 4. B & 5. C & 6. A & 7. B & 8. C & 9. B & 10. D\\\hline
11. D & 12. D & 13. D & 14. B & 15. A & 16. A & 17. D & 18. C & 19. A & 20. D\\\hline
21. D & 22. D & 23. D & 24. B & 25. D & 26. D & 27. C & 28. D & 29. D & 30. B\\\hline
31. C & 32. D & 33. D & 34. B & 35. A & 36. C & 37. C & 38. B & 39. C & 40. C\\\hline
41. A & 42. A & 43. C & 44. A & 45. D & 46. B & 47. C & 48. A & 49. C & 50. C\\\hline
\end{tabular}\end{center}
\end{document}
